\chapter{Machine Learning}\label{ch:iv:machine-learning}

A candidate's model scored 71\% accuracy in cross-validation on the daily direction of a futures contract and 50\% in its
first month live. Asked why, he says ``overfitting''. The interviewer asks him to name the three lines of his pipeline most
likely to have leaked the future, and the interview starts there. Machine-learning questions in finance interviews are
rarely about architectures. They are about the ways a model can look good and be worthless, which in markets are more
numerous than in most fields: the signal is weak, the data are few and dependent, and the future does not resemble the
past. The vocabulary is One Quant Book~12's (chapters 1 to 11) and the validation discipline One Quant Book~7's
(chapters 3 and~20).

\section{Bias, variance and regularisation}

The bias--variance decomposition (One Quant Book~12, chapter~1) answers most ``what happens if'' questions. More data lowers
variance and leaves bias; more features or a more flexible model lower bias and raise variance; more regularisation raises
bias and lowers variance. In finance, where the signal-to-noise ratio of returns is tiny, variance dominates: the best
models are usually simple, strongly regularised, and evaluated honestly.

Ridge and lasso regularise differently (One Quant Book~4, chapter~16). The ridge penalty $\lambda\|\beta\|_2^2$ shrinks all
coefficients and shares weight among correlated features; the lasso penalty $\lambda\|\beta\|_1$ has corners on the axes, so
its solution often sits on one, setting some coefficients exactly to zero. With nearly duplicate features the lasso picks
one of them, and which one depends on noise.

\begin{example}[Two nearly identical features]\label{ex:iv:machine-learning:lasso}
Two features correlated at 0.99, a target equal to their sum plus noise, 100 observations. Over 200 samples, a lasso with
$\alpha = 0.5$ sets one of the two coefficients to zero in 63\% of them, choosing each about equally often; ridge gives both
about 0.95 every time. The lasso's selection is not a finding about which feature matters.
\end{example}

\section{Trees, ensembles and networks}

Gradient boosting and random forests (One Quant Book~12, chapter~5) fit interactions and non-linearities without
specification, which is why they overfit noise readily: with enough trees a boosted model fits a pure-noise training set.

\begin{example}[Boosting on noise]\label{ex:iv:machine-learning:boost}
Five noise features and a noise target, 1\,000 training and 1\,000 test observations, 300 trees of depth 3 at a learning
rate of 0.1: training $R^2$ of 0.65 and test $R^2$ of $-0.18$. Early stopping on a validation set that respects time would
have stopped near zero trees, which is the right answer.
\end{example}

Two further traps recur. Impurity-based feature importances favour continuous and high-cardinality features even when they
are noise; use permutation importance on held-out data. And deep networks on daily returns disappoint for reasons of data,
not architecture: a few thousand dependent observations of a signal with an out-of-sample $R^2$ well below 1\% cannot
train millions of parameters (One Quant Book~12, chapter~7).

\section{Validation with time}

\begin{method}[Finding the leak]\label{met:iv:machine-learning:leak}
\begin{enumerate}
\item \emph{Timestamps}: is every feature computed from data available at the decision time, with its publication lag?
(point-in-time data, One Quant Book~7, chapter~3)
\item \emph{Labels}: do labels overlap in time (a 20-day forward return every day)? Then neighbouring observations share
most of their label.
\item \emph{Splits}: does cross-validation shuffle observations, so that neighbours of a test point sit in the training set?
Purge training observations whose labels overlap the test period and add an embargo after it (One Quant Book~7,
chapter~20).
\item \emph{Preprocessing}: were scalers, feature selection or hyperparameters fitted on the whole sample?
\item \emph{Selection}: how many models were tried before this one?
\end{enumerate}
\end{method}

\begin{omfigure}
\begin{tikzpicture}
\begin{axis}[width=10cm,height=5.2cm,xlabel={\small simulated market (seed)},ylabel={\small out-of-sample $R^2$},
  xmin=-0.5,xmax=19.5,ymin=-1.1,ymax=0.5,xtick={0,5,10,15,19},
  tick label style={font=\footnotesize},grid=major,grid style={gray!20},table/col sep=comma,
  legend cell align=left,legend style={font=\scriptsize,at={(0.5,-0.32)},anchor=north,
  /tikz/every even column/.append style={column sep=8pt}},legend columns=2]
\addplot[only marks,mark=*,mark size=1.8pt,omThm] table[x=seed,y=shuffled]{figdata/interviews/19-machine-learning/leak.csv};
\addplot[only marks,mark=square*,mark size=1.6pt,omDef] table[x=seed,y=walkforward]{figdata/interviews/19-machine-learning/leak.csv};
\draw[gray] (axis cs:-0.5,0) -- (axis cs:19.5,0);
\legend{shuffled 5-fold,purged walk-forward}
\end{axis}
\end{tikzpicture}

\omcaption{Twenty simulated random walks with no predictability, trailing-return features and overlapping 20-day forward
labels, fitted with five-nearest-neighbour regression. Shuffled cross-validation reports a positive $R^2$ on every one
(mean 0.26); a walk-forward evaluation with purging reports a negative one (mean $-0.61$). Data:
\texttt{fig\_iv\_leak.py}.}\label{fig:iv:machine-learning:leak}
\end{omfigure}

The mechanism in \cref{fig:iv:machine-learning:leak} is leakage through overlap: slowly varying features make yesterday and
tomorrow the nearest neighbours of today, and overlapping labels make their targets almost today's target, so a shuffled
split lets the model look up the answer. Nothing is predictable; the score is.

\section{Design questions}

\begin{definition}[Machine-learning design question]\label{def:iv:machine-learning:design}
A \emph{machine-learning design question}\index{machine-learning design question} asks the candidate to turn a vague business
goal into a learning problem: the decision the model serves, the target and its horizon, the features available at decision
time, the validation scheme, the metric that reflects the decision's payoff, and the monitoring after deployment.
\end{definition}

\begin{method}[Answering a design question]\label{met:iv:machine-learning:design}
\begin{enumerate}
\item Name the decision and its cost structure first; the metric follows from it, not from convention.
\item Define the target precisely (event, horizon, how it is observed, when it is known).
\item List features by source and latency; exclude anything not known at the decision time.
\item Start with a baseline (a rule, a logistic regression), then a stronger model if it earns its complexity.
\item Validate forward in time with purging; report the metric with its uncertainty.
\item Say how the model fails in production (drift, feedback from its own actions) and what is monitored
(One Quant Book~12, chapter~27).
\end{enumerate}
\end{method}

\section{Worked answers}

\begin{example}[Reading a confusion table aloud]\label{ex:iv:machine-learning:confusion}
\emph{``Of 10\,000 orders, 3\% come from counterparties later judged toxic. The model flags 500 orders, of which 240 are
toxic. How good is it?''} There are 300 toxic orders; the model catches 240, a recall of 80\%, and 240 of its 500 flags are
right, a precision of 48\%. Against the base rate of 3\%, a flagged order is 16 times more likely to be toxic than a random
one: the lift. Whether 48\% is good depends on the action: widening quotes for the flagged orders costs a little on the 260
false alarms and saves a lot on the 240 toxic ones, and the threshold should be moved until the marginal flag's expected
saving equals its expected cost (\cref{iq:iv:machine-learning:13} does this for a trade).
\end{example}

\begin{example}[How many independent labels?]\label{ex:iv:machine-learning:overlap}
\emph{``We have ten years of daily data and predict twenty-day forward returns. How many observations do we have?''}
About 2\,520 rows but far fewer independent labels: consecutive twenty-day returns share nineteen days, so roughly
$2\,520/20 = 126$ non-overlapping ones. A $t$-statistic computed as if the 2\,520 were independent is inflated by about
$\sqrt{20} \approx 4.5$, and a validation fold whose labels overlap the training fold's leaks. The fixes are standard:
purge training rows whose label window overlaps the validation period, embargo a gap after it, and compute standard errors
with overlap-robust (Newey--West or block bootstrap) methods (\cref{met:iv:machine-learning:leak}).
\end{example}

\section{Question bank}

\begin{interviewq}[$\star$ \iqroles{mle, researcher} \iqfirm{systematic fund}]\label{iq:iv:machine-learning:1}
What happens to bias and to variance when you add data, add features, or increase regularisation?
\end{interviewq}

\begin{interviewq}[$\star$ \iqroles{mle, researcher} \iqfirm{any}]\label{iq:iv:machine-learning:2}
Why does the lasso set coefficients exactly to zero while ridge does not?
\end{interviewq}

\begin{interviewq}[$\star$ \iqroles{researcher, trader} \iqfirm{market maker}]\label{iq:iv:machine-learning:3}
A model predicts the next day's direction with 55\% accuracy. Is that good?
\end{interviewq}

\begin{interviewq}[$\star$ \iqroles{mle} \iqfirm{any}]\label{iq:iv:machine-learning:4}
Why standardise features before fitting ridge or lasso, and with which statistics?
\end{interviewq}

\begin{interviewq}[$\star\star$ \iqroles{mle, researcher} \iqfirm{systematic fund}]\label{iq:iv:machine-learning:5}
Two features are correlated at 0.99. A lasso keeps one and drops the other, and on a new sample it keeps the other one.
What is happening, and what do you report?
\end{interviewq}

\begin{interviewq}[$\star\star$ \iqroles{researcher, mle} \iqfirm{systematic fund}]\label{iq:iv:machine-learning:6}
A colleague predicts 20-day forward returns every day from trailing 20-, 60- and 120-day returns, with shuffled 5-fold
cross-validation, and reports an $R^2$ of 0.26. What do you suspect, and how do you test it?
\end{interviewq}

\begin{interviewq}[$\star\star$ \iqroles{mle} \iqfirm{proprietary firm}]\label{iq:iv:machine-learning:7}
A gradient-boosted model has a training $R^2$ of 0.65 and a test $R^2$ of $-0.18$. What happened, and what would you change?
\end{interviewq}

\begin{interviewq}[$\star\star$ \iqroles{researcher, mle} \iqfirm{systematic fund}]\label{iq:iv:machine-learning:8}
A random forest's feature importances rank a noisy continuous feature first. Why can that happen, and what importance would
you use instead?
\end{interviewq}

\begin{interviewq}[$\star\star$ \iqroles{mle, researcher} \iqfirm{multi-manager fund}]\label{iq:iv:machine-learning:9}
Why does a deep network trained on ten years of daily returns usually do worse than a regularised linear model?
\end{interviewq}

\begin{interviewq}[$\star\star\star$ \iqroles{mle, developer} \iqfirm{market maker}]\label{iq:iv:machine-learning:10}
Design a model that predicts the probability that a passive limit order will be filled within one second.
\end{interviewq}

\begin{interviewq}[$\star\star\star$ \iqroles{mle, trader} \iqfirm{market maker}]\label{iq:iv:machine-learning:11}
Design a model that flags toxic order flow for a market maker, so that quotes can be widened for some counterparties.
\end{interviewq}

\begin{interviewq}[$\star\star\star$ \iqroles{researcher, mle} \iqfirm{systematic fund}]\label{iq:iv:machine-learning:12}
A model scored 71\% accuracy in cross-validation on daily direction and 50\% in its first month live. List the most likely
causes in order, and the check for each.
\end{interviewq}

\begin{interviewq}[$\star\star\star$ \iqroles{researcher, trader} \iqfirm{proprietary firm}]\label{iq:iv:machine-learning:13}
A classifier outputs the probability that a trade idea is right. A right trade gains 1 and a wrong one loses 2. Above what
probability do you trade? Why is accuracy the wrong metric, and which would you use?
\end{interviewq}

\begin{omsources}
\item One Quant Book~12, chapters 1--11 and~27; One Quant Book~7, chapters 3 and~20; One Quant Book~4, chapter~16.
\item M.~L\'opez de Prado, \emph{Advances in Financial Machine Learning}, Wiley, 2018, on purging and embargoes.
\end{omsources}
